\documentclass[a4paper]{article}

\usepackage[a4paper, top=20mm, bottom=20mm, left=20mm, right=20mm]{geometry}

\usepackage{polyglossia}
\setdefaultlanguage[babelshorthands=true]{russian}
\setotherlanguage{english}

\usepackage{fontspec}
\setmainfont{FreeSerif}
\newfontfamily{\russianfonttt}[Scale=0.7]{DejaVuSansMono}

\usepackage{hyperref}
\usepackage{bookmark}
\usepackage{csquotes}

\title{Обзор алгоритма Хаффмана и существующих реализаций}
\author{Роман Казачков}
\date{}

\begin{document}

\maketitle

\section{Обзор алгоритма и существующих реализаций}

Алгоритмы сжатия данных без потерь применяются в тех случаях, когда исходная
информация должна быть восстановлена полностью. Такие методы используются в
архиваторах, системах хранения данных, сетевых протоколах и различных форматах
файлов. Одним из наиболее известных алгоритмов такого типа является кодирование
Хаффмана.

Алгоритм был предложен Дэвидом Хаффманом в 1952 году в работе \enquote{A Method
for the Construction of Minimum-Redundancy Codes}\footnote{D. A. Huffman.
\enquote{A Method for the Construction of Minimum-Redundancy Codes}.
\url{https://ieeexplore.ieee.org/document/4051119}.}. Он относится к методам энтропийного кодирования
и основан на различной частоте появления символов в данных. Часто встречающиеся
символы кодируются короткими битовыми последовательностями, а редкие --- более
длинными. Это позволяет уменьшить среднюю длину кодового слова и сократить
общий размер файла.

Работа алгоритма начинается с подсчёта частоты появления символов во входных
данных. На основе этой статистики строится бинарное дерево. На каждом шаге
выбираются два узла с минимальной частотой и объединяются в новый внутренний
узел. Процесс продолжается до тех пор, пока не останется один корневой узел.
После этого выполняется обход дерева для генерации кодов: переход влево
соответствует биту 0, переход вправо --- биту 1. Полученные коды обладают
свойством префиксности, поэтому их можно однозначно декодировать.

Алгоритм Хаффмана используется как часть более сложных методов сжатия.
Например, он применяется в алгоритме \enquote{DEFLATE}, который лежит в основе
форматов ZIP и gzip. Реализация этого алгоритма присутствует в библиотеке
zlib\footnote{Библиотека zlib. \url{https://zlib.net}.}. Кодирование Хаффмана также используется в
некоторых форматах изображений и мультимедийных данных, например в JPEG.

Существует несколько вариантов алгоритма. В статической версии дерево строится
один раз на основе статистики всего файла. В адаптивных версиях структура
дерева может изменяться в процессе обработки данных, что позволяет работать с
потоками информации без предварительного анализа.

Алгоритм Хаффмана реализован во многих библиотеках и учебных проектах. Примеры
открытых реализаций можно найти в исходном коде zlib\footnote{Исходный код
zlib. \url{https://github.com/madler/zlib}.}, а также в демонстрационных проектах на GitHub:
реализации на языке C\footnote{Реализация алгоритма Хаффмана на языке C.
\url{https://github.com/drichardson/huffman}.} и C++\footnote{Реализация алгоритма Хаффмана на
языке C++. \url{https://github.com/TheAlgorithms/C-Plus-Plus/blob/master/greedy_algorithms/huffman.cpp}.}.

Большое количество доступных реализаций делает алгоритм удобным для изучения и
экспериментов. Простота идеи и эффективность привели к тому, что кодирование
Хаффмана до сих пор используется как компонент более сложных систем сжатия
данных.

\end{document}
